\documentclass{article}
\usepackage{graphicx, indentfirst, amssymb, amsmath, amsthm, placeins, xurl, hyperref, multicol, algorithm, algorithmic} 

\title{Bandit Convex Optimization, Chap 6.1 - 6.3, Presentation Notes}
\author{Artem Ivaniuk}
\date{\today}

\begin{document}
\maketitle

\section{BCO Intro}
In real-world applications of Online Convexity Optimization, the feedback / loss function revealed to the player is incomplete.

For Bandit Convex Optimization (BCO), we take the standard OCO setting and instead of the player finding out the entire gradient oracle (loss function) $f_t$, the player only has access to $f_t(x_t)$ where $x_t$ is the action played at time $t$.

Set up:
\begin{enumerate}
\item At iteration $t$, the online player chooses $x_t \in \mathcal K$. 
\item Then, a convex cost function $f_t \in \mathcal{F}: \mathcal{K} \to \mathbb{R}$ is chosen (not revealed to the player). $\mathcal F$ being the total set of functions available to the adversary
\item The cost/loss incurred to the online player is $f_t(x_t)$ 
\item \textbf{Main change}: only loss $f_t(x_t)$ is revealed to the player, not entire $f_t$
\item For algorithm $\mathcal A$, same Regret definition:
\[\operatorname{Regret}_T(\mathcal A) = \operatorname{sup}_{\{f_1, \dots, f_T\}\subseteq \mathcal F} \left\{\sum_{t=1}^T f_t(x_t) - \min_{x \in \mathcal K} \sum_{t=1}^T f_t(x)\right\}\]
\end{enumerate}

\section{Multi-Armed Bandit Reduction}

Let $\ell_t \in \mathbb{R}^n$ be a loss vector, where $\ell_t(i)$ is the loss of action $i$.

We define a linear function:
\[
f_t(x) = \ell_t^\top x
\]

If we choose action $i_t$, this corresponds to evaluating:
\[
f_t(i_t) = \ell_t(i_t)
\]

So the bandit setting is a special case of OCO with linear functions.

The main issue is that we only observe $\ell_t(i_t)$ but we need information about all coordinates of $\ell_t$

\section{Simple MAB Algorithm (Algorithm 20)}

\begin{algorithm}
\caption{Simple MAB algorithm}
\begin{algorithmic}[1]
\REQUIRE OCO algorithm $A$, parameter $\delta$
\FOR{$t = 1$ to $T$}
    \STATE Let $b_t$ be a Bernoulli random variable that equals $1$ with probability $\delta$
    \IF{$b_t = 1$}
        \STATE Choose $i_t \in \{1,2,\dots,n\}$ uniformly at random and play it
        \STATE Let
        $
        \hat{\ell}_t(i) =
        \begin{cases}
        \frac{n}{\delta}\,\ell_t(i_t), & i = i_t \\
        0, & \text{otherwise}
        \end{cases}
        $
        \STATE Let $\hat{f}_t(x) = \hat{\ell}_t^\top x$ and update $x_{t+1} = A(\hat{f}_1,\dots,\hat{f}_t)$
    \ELSE
        \STATE Choose $i_t$ according to $x_t$ and play it
        \STATE Set $\hat{\ell}_t = 0$, $\hat{f}_t = 0$, and $x_{t+1} = x_t$
    \ENDIF
\ENDFOR
\end{algorithmic}
\end{algorithm}

Estimator:
\[
\hat{\ell}_t(i) =
\begin{cases}
\frac{n}{\delta} \ell_t(i_t) & i = i_t \\
0 & \text{otherwise}
\end{cases}
\]

We only observe $\ell_t(i_t)$ when we explore and pick $i_t$ uniformly. This happens with probability $\delta \cdot \frac{1}{n}$. To compensate for this, we divide by that probability:
\[
\frac{1}{\delta \cdot (1/n)} = \frac{n}{\delta}
\]

This gives:
\[
\mathbb{E}[\hat{\ell}_t(i)] = \ell_t(i)
\]

So the estimator is unbiased. The idea is that we rarely see each coordinate, so when we do see it, we scale it up to make it representative.

\section{Lemma 6.2}

We want to show:
\[
\mathbb{E}[\ell_t(i_t)] \le \mathbb{E}[\hat{\ell}_t^\top x_t] + \delta
\]

Split the expectation based on whether we explore or exploit:
\[
\mathbb{E}[\ell_t(i_t)] = \delta \mathbb{E}[\ell_t(i_t)\mid b_t=1] + (1-\delta)\mathbb{E}[\ell_t(i_t)\mid b_t=0]
\]

In the exploration case, losses are bounded in $[0,1]$, so:
\[
\mathbb{E}[\ell_t(i_t)\mid b_t=1] \le 1
\]

In the exploitation case, $i_t$ is sampled from $x_t$, so:
\[
\mathbb{E}[\ell_t(i_t)\mid b_t=0] = \ell_t^\top x_t
\]

Combining:
\[
\mathbb{E}[\ell_t(i_t)] \le \delta + \ell_t^\top x_t
\]

Since $\mathbb{E}[\hat{\ell}_t] = \ell_t$, we have:
\[
\mathbb{E}[\ell_t^\top x_t] = \mathbb{E}[\hat{\ell}_t^\top x_t]
\]

So:
\[
\mathbb{E}[\ell_t(i_t)] \le \mathbb{E}[\hat{\ell}_t^\top x_t] + \delta
\]

The idea is that exploitation behaves like the expected loss under $x_t$, and exploration adds at most $\delta$ error.

\section{Lemma 6.1}

We want to bound regret.

We decompose:
\[
\text{Regret} = \text{regret on } \hat{\ell}_t + \delta T
\]

The first term comes from running OCO on $\hat{\ell}_t$. The second term comes from exploration, since we explore with probability $\delta$ each round.

The estimator can be large:
\[
\|\hat{\ell}_t\| \le \frac{n}{\delta}
\]

This happens because we divide by a small probability. So rare observations get scaled up significantly.

The regret behaves like:
\[
\frac{n}{\delta}\sqrt{\delta T} + \delta T
\]

First term comes from variance in the estimator, second from exploration.

We balance these two terms by choosing:
$\delta = n^{2/3} T^{-1/3}$

This gives regret:
$O(T^{2/3} n^{2/3})$

The intuition is that we trade off between exploring enough to learn and not exploring too much to avoid large loss

\pagebreak
\section{EXP3 Algorithm (Algorithm 21)}
\FloatBarrier
\begin{algorithm}
\caption{EXP3 - simple version}
\begin{algorithmic}[1]
\REQUIRE parameter $\varepsilon > 0$
\STATE Set $x_1 = \left(\frac{1}{n},\dots,\frac{1}{n}\right)$
\FOR{$t = 1$ to $T$}
    \STATE Choose $i_t$ according to $x_t$ and play it
    \STATE Let
    $
    \hat{\ell}_t(i) =
    \begin{cases}
    \frac{1}{x_t(i_t)}\,\ell_t(i_t), & i = i_t \\
    0, & \text{otherwise}
    \end{cases}
    $
    \STATE Update $y_{t+1}(i) = x_t(i)\exp(-\varepsilon \hat{\ell}_t(i))$
    \STATE Set $x_{t+1} = \frac{y_{t+1}}{\|y_{t+1}\|_1}$
\ENDFOR
\end{algorithmic}
\end{algorithm}
\FloatBarrier

Estimator:
\[
\hat{\ell}_t(i) =
\begin{cases}
\frac{\ell_t(i_t)}{x_t(i_t)} & i = i_t \\
0 & \text{otherwise}
\end{cases}
\]

We observe coordinate $i$ with probability $x_t(i)$, so we divide by $x_t(i)$ to correct for this. This again ensures:
\[
\mathbb{E}[\hat{\ell}_t(i)] = \ell_t(i)
\]

\section{Lemma 6.3}

We bound:
\[
\mathbb{E}[\text{Regret}_T] \le 2\sqrt{T n \log n}
\]

We analyze:
\[
\mathbb{E}[x_t^\top \hat{\ell}_t^2]
\]

Only one coordinate of $\hat{\ell}_t$ is nonzero, so:
\[
x_t(i_t)\left(\frac{\ell_t(i_t)}{x_t(i_t)}\right)^2 = \frac{\ell_t(i_t)^2}{x_t(i_t)}
\]

Taking expectation:
\[
\sum_i x_t(i)\frac{\ell_t(i)^2}{x_t(i)} = \sum_i \ell_t(i)^2 \le n
\]

So although $\hat{\ell}_t$ can be large, it is large with small probability, and the expectation remains bounded.

Using the Hedge bound:
\[
\text{Regret} \le \varepsilon T n + \frac{\log n}{\varepsilon}
\]

Choosing:
\[
\varepsilon = \sqrt{\frac{\log n}{T n}}
\]

gives the result.

The improvement over the simple algorithm is that exploration happens every round, but in a controlled way, reducing variance.

\section{General Reduction (Algorithm 22)}

\begin{algorithm}
\caption{Reduction to bandit feedback}
\begin{algorithmic}[1]
\REQUIRE convex set $K \subseteq \mathbb{R}^n$, first order online algorithm $A$
\STATE Let $x_1 = A(\emptyset)$
\FOR{$t = 1$ to $T$}
    \STATE Generate distribution $D_t$, sample $y_t \sim D_t$ with $\mathbb{E}[y_t] = x_t$
    \STATE Play $y_t$
    \STATE Observe $f_t(y_t)$, generate $g_t$ with $\mathbb{E}[g_t] = \nabla f_t(x_t)$
    \STATE Let $x_{t+1} = A(g_1,\dots,g_t)$
\ENDFOR
\end{algorithmic}
\end{algorithm}

The idea is to simulate gradients using an estimator $g_t$ that satisfies:
\[
\mathbb{E}[g_t] = \nabla f_t(x_t)
\]

\section{Lemma 6.5}

Define:
\[
\xi_t = g_t - \nabla f_t(x_t)
\]

and:
\[
h_t(x) = f_t(x) + \xi_t^\top x
\]

Then:
\[
\nabla h_t(x_t) = g_t
\]

So the algorithm is effectively running OCO on $h_t$.

Applying the OCO guarantee:
\[
\sum h_t(x_t) - \sum h_t(u) \le B_A
\]

Expanding:
\[
\sum f_t(x_t) + \sum \xi_t^\top x_t - \sum f_t(u) - \sum \xi_t^\top u
\]

Since:
\[
\mathbb{E}[\xi_t] = 0
\]

the noise terms vanish in expectation, giving:
\[
\mathbb{E}\left[\sum f_t(x_t)\right] - \sum f_t(u) \le \mathbb{E}[B_A]
\]

The idea is that we hide the noise inside the function, and expectation removes it.

\section{Gradient Estimation}

In one dimension:
\[
f'(x) \approx \frac{f(x+\varepsilon) - f(x-\varepsilon)}{2\varepsilon}
\]

But in the bandit setting, we only get one query.

So we use randomness. In higher dimensions:
\[
g = \frac{n}{\varepsilon} f(x+\varepsilon u) u
\]

where $u$ is a random direction.

Then:
\[
\mathbb{E}[g] = \nabla f(x)
\]

Each sample gives a directional estimate of the slope, and averaging over directions reconstructs the full gradient.

\end{document}